% Short note: projection-input recognition + decoupled readout estimation for VJF.
% Notation follows Zhao & Park, "Variational online learning of neural dynamics"
% (Frontiers Comput. Neurosci. 2020); eq. labels reference that manuscript.
\documentclass[11pt]{article}
\usepackage[margin=2cm]{geometry}
\usepackage{amsmath,amssymb,bm}
\newcommand{\vx}{\bm{x}}\newcommand{\vy}{\bm{y}}\newcommand{\vu}{\bm{u}}
\newcommand{\vC}{\bm{C}}\newcommand{\vb}{\bm{b}}\newcommand{\vW}{\bm{W}}
\newcommand{\vR}{\bm{R}}\newcommand{\vmu}{\bm{\mu}}\newcommand{\vs}{\bm{s}}
\newcommand{\vphi}{\bm{\phi}}\newcommand{\vl}{\bm{\ell}}\newcommand{\vnu}{\bm{\nu}}
\newcommand{\vpi}{\bm{\pi}}\newcommand{\E}{\mathbb{E}}\newcommand{\R}{\mathbb{R}}
\newcommand{\vc}{\bm{c}}
\newcommand{\half}{\tfrac12}
\begin{document}
\begin{center}\large\textbf{Projection-input recognition and decoupled readout estimation}\end{center}

\paragraph{Model recap (unchanged).}
With latent $\vx_t\in\R^m$, observation $\vy_t\in\R^n$, input $\vu_t$, the generative
model is the point-process / GLM observation and additive RBF dynamics of the paper,
\begin{align}
  \vy_t &\sim P\big(g(\vC\vx_t+\vb)\big), \qquad g=\exp \tag{lik}\\
  \vx_{t+1} &= \vx_t + f(\vx_t) + \bm{B}_t\vu_t + \bm{\epsilon}_{t+1},\quad
  f(\vx)=\vW\vphi(\vx),\quad \bm{\epsilon}_t\sim\mathcal N(\bm0,\sigma^2\bm I), \tag{gen}
\end{align}
and the recursive amortized posterior $q(\vx_t)=\mathcal N(\vmu_t,\vs_t)$ is produced by a
recognition MLP $h$, optimizing the reparameterized filtering ELBO
\begin{align}
  [\vmu_t,\vs_t] &= h(\vy_t,\vu_{t-1},\vmu_{t-1},\vs_{t-1}), \tag{recog}\\
  \hat{\mathcal L} &= \log p(\vy_t\mid\tilde\vx_t,\theta)
     + \E_{q(\vx_t)}\log p(\vx_t\mid\tilde\vx_{t-1},\theta) + H(q(\vx_t)). \tag{ELBO}
\end{align}
The paper notes (lik) is identifiable only up to an invertible $\vR$:
$\vC\vx_t=(\vC\vR)(\vR^{-1}\vx_t)$, and ``normalizes $\vC$ each iteration.''

\paragraph{The problem.}
Estimating $\vC$ by $\partial\hat{\mathcal L}/\partial\vC$ under sparse, low-rate Poisson
collapses to the trivial ``mean rate via $\vb$'' solution: the gradient on $\vC$ is scaled
by the rate $g(\vC\vx_t+\vb)\!\approx\!10^{-1}$/bin, so $\vC$ barely moves and a basin with
$\vC\!\to\!0$ absorbs the loss. Separately, $h(\vy_t,\dots)$ must invert $n$ sparse channels,
an $n\!\to\!m$ map that is hard to learn at low SNR.

\paragraph{Update (two coupled changes; the ELBO and generative model are unchanged).}
\textbf{(1) Recognition input.} Replace the raw observation $\vy_t$ in (recog) by a
deterministic, $\vC$-dependent feature --- the least-squares projection of a link-matched
observation feature onto the current loading subspace:
\begin{align}
  \vpi_t &= \vC^{+}\,(\tilde g(\vy_t)-\vb), \qquad \vC^{+}=(\vC^\top\vC)^{-1}\vC^\top\in\R^{m\times n},
     &&\text{(subspace projection)}\\
  [\vmu_t,\vs_t] &= h(\vpi_t,\vu_{t-1},\vmu_{t-1},\vs_{t-1}), &&\text{(recog$'$)}
\end{align}
where $\tilde g(\cdot)$ is an observation-model-specific, variance-stabilizing feature
(an approximate inverse link) and $\vb$ is the running mean of $\tilde g(\vy_t)$. For the
point-process/Poisson case $\tilde g(\vy_t)=\log(\vnu_t+c)$ with a causally-smoothed count
$\vnu_t=(1-\alpha)\vnu_{t-1}+\alpha\vy_t$ ($\alpha=1/\tau$); for Gaussian observations
$\tilde g=$ identity.
$\vpi_t$ is the readout's own approximate inverse, so $h$ need only perform an $m\!\to\!m$
correction rather than invert $n$ channels. Because $q(\vx_t)$ may be \emph{any}
data-dependent distribution, parameterizing it through $\vpi_t$ leaves the bound
$\hat{\mathcal L}$ valid (same likelihood, dynamics, entropy); only the amortization family
changes. $\vC^{+}$ is recomputed only when $\vC$ is refreshed (below) and reused as one
$O(nm)$ map per step --- not re-solved per bin. Any feature timescale $\tau$ is a fixed
hyperparameter ($\tau\ll$ the dynamics timescale), \emph{not} ELBO-learned: it affects only
the recognition input ($h$ also sees $\vmu_{t-1}$ to compensate the lag), is differentiable
in $\alpha$, and could be learned, but we set or sweep it.

\textbf{(2) Decoupled online readout estimation.} Replace $\partial\hat{\mathcal L}/\partial\vC$
by an online incremental-PCA estimate $\hat\vC_t$ of the centered smoothed log-rate: track
the top-$m$ subspace of $\vl_t-\bar\vb_t$ (e.g. CCIPCA/Oja), scaled so the implied latent has
unit variance, with $\vb=\bar\vb_t$ the running mean of $\vl_t$. The \emph{same} $\hat\vC_t$
appears in the decoder (lik) and in the projection $\vpi_t$. This is an online M-step
(second-moment estimator), well-conditioned and free of the rate-scaled-gradient collapse.
$\hat\vC_t$ is refreshed every $K$ steps.

\textbf{(3) Identifiability / normalization.} The paper's ``normalize $\vC$ each iteration''
(for $\vC\vx_t=(\vC\vR)(\vR^{-1}\vx_t)$) becomes: on each refresh choose $\vR$ by orthogonal
Procrustes to align $\hat\vC_t$ to the previous estimate, fixing the latent coordinate frame
so the dynamics $f=\vW\vphi$ stay valid. For a pure rotation $\vR$ the RBF flow is
re-expressed exactly by rotating the centers $\vc_i\mapsto\vR\vc_i$ (and $\vW$ accordingly).

\paragraph{Special cases.}
$\hat\vC_t\equiv\vC$ (known/oracle readout) recovers a fixed projection; $\vpi_t:=\vy_t$ with
$\vC$ learned by $\partial\hat{\mathcal L}/\partial\vC$ recovers the original VJF. Per-step
cost stays $O(nm)$ (same order as the decoder), preserving real-time complexity.

\paragraph{Behavior.}
From a cheap small-window initial $\hat\vC$, the decoupled estimator recovers the asymptotic
readout subspace online, and the projection-input recognition attains the known-readout
(oracle) filtering accuracy --- and can \emph{exceed} raw-$\vy_t$ recognition, since $h$ no
longer has to invert $n$ (possibly sparse) channels.
\end{document}
